\section{Method}
\label{sec:method}

\begin{figure}[t]
\centering
\begin{tikzpicture}[
  font=\scriptsize,
  stage/.style={draw=black!55, rounded corners=2pt, align=center, inner sep=2.5pt, text width=2.05cm, minimum height=1.05cm, line width=0.5pt, fill=black!3},
  ours/.style={stage, draw=cvprblue, fill=cvprblue!10, line width=0.8pt},
  arr/.style={-{Latex[length=1.5mm,width=1.2mm]}, line width=0.6pt, black!70},
  darr/.style={-{Latex[length=1.3mm,width=1mm]}, line width=0.5pt, black!55, dashed},
  pt/.style={circle, fill=black!75, inner sep=1.2pt},
  apt/.style={circle, draw=black!60, fill=white, inner sep=1.2pt, line width=0.5pt},
  lab/.style={font=\tiny, black!75, inner sep=1pt},
]
% ---- pipeline row -------------------------------------------------------
\node[stage] (mid) at (0,0) {\textbf{Mid-training}\\ web video +\\ cross-embodiment\\ robot data\\[1pt] $\Rightarrow\theta_{\mathrm{mid}}$};
\node[ours, right=0.38 of mid] (meta) {\textbf{\Method}\\ \textbf{(ours)}\\ 5 source embodiments,\\ no new data\\[1pt] $\theta_{\mathrm{mid}}\Rightarrow\thmeta$};
\node[stage, right=0.38 of meta] (post) {\textbf{Post-training}\\ target robot,\\ 3 demos / task,\\ unchanged recipe\\[1pt] $\thmeta\Rightarrow$ policy};
\draw[arr] (mid) -- (meta);
\draw[arr] (meta) -- (post);

% ---- zoom box: one meta-iteration -------------------------------------
\path let \p1=(mid.west), \p2=(post.east) in node[draw=cvprblue!70, rounded corners=2pt, line width=0.5pt, fill=cvprblue!3, anchor=north west, inner sep=0pt, minimum width={\x2-\x1}, minimum height=2.8cm] (zoom) at ($(mid.south west)+(0,-0.42)$) {};
\draw[cvprblue!70, line width=0.5pt] (meta.south west) -- ($(zoom.north west)!0.28!(zoom.north east)$);
\draw[cvprblue!70, line width=0.5pt] (meta.south east) -- ($(zoom.north west)!0.72!(zoom.north east)$);
\node[lab, cvprblue, anchor=north west, font=\tiny\bfseries] at ($(zoom.north west)+(0.08,-0.05)$) {One meta-iteration};

% parameter-space sketch (left)
\begin{scope}[shift={($(zoom.west)+(0.75,-0.12)$)}]
  \node[pt, label={[lab]left:$\theta$}] (th) at (0,0) {};
  \node[apt, label={[lab]above right:YAM}] (a5) at (0.15,0.72) {};
  \node[apt, label={[lab]above right:RT-1}] (a1) at (1.40,0.66) {};
  \node[apt, label={[lab]right:Bridge}] (a2) at (1.85,0.0) {};
  \node[apt, label={[lab]below:UR}] (a3) at (1.40,-0.70) {};
  \node[apt, label={[lab]below:Franka}] (a4) at (0.35,-0.82) {};
  \foreach \a in {a1,a2,a3,a4,a5} \draw[darr] (th) -- (\a);
  \draw[arr, cvprblue, line width=0.9pt] (th) -- ($(th)!0.45!(a2)$);
  \node[lab] at ($(th)!0.58!(a1)+(0.02,0.15)$) {$\tilde\theta_e$};
  \node[lab, cvprblue, anchor=north] at ($(th)!0.3!(a2)+(0.05,-0.06)$) {$\epsilon(\tilde\theta_e-\theta)$};
\end{scope}

% recipe text (right)
\node[anchor=north west, align=left, text width=3.8cm, inner sep=2pt] at ($(zoom.north west)+(3.45,-0.3)$) {%
  1.\ sample embodiment $e$, $K{=}32$ demos, $w{=}8$ windows each\\[1pt]
  2.\ $\tilde\theta_e\leftarrow$ $k{=}10$ steps of the \emph{post-training} optimizer on them (same loss, learning rate, trainable set; shared action head)\\[1pt]
  3.\ $\theta\leftarrow\theta+\epsilon\,(\tilde\theta_e-\theta)$ with \mbox{$\epsilon\colon 0.5\to0.05$}%
};
\end{tikzpicture}
\caption{\textbf{Overview.}
We insert a \method stage between mid-training and post-training.
Each meta-iteration simulates few-shot post-training on one source embodiment ($K$ demonstrations, $k$ steps of the unchanged post-training recipe; dashed arrows to the adapted weights $\tilde\theta_e$) and moves the weights a fraction $\epsilon$ toward the adapted solution (blue arrow).
The result $\thmeta$ is loaded by the post-training recipe exactly like the mid-trained checkpoint.}
\label{fig:method}
\end{figure}


We describe our method with \base~\cite{cosmos3} as the representative backbone.
The method only assumes a pre-trained policy backbone trained with an action objective, a post-training recipe, and a set of source robot datasets, and applies to VLAs in the same way.
We first describe the backbone and its post-training recipe (\cref{sec:prelim}), then define the few-shot post-training problem (\cref{sec:problem}), present \Method (\cref{sec:meta}) and \Rpl (\cref{sec:rpl}), and finally summarize how post-training uses them (\cref{sec:downstream}).
\cref{fig:method} gives an overview.

\subsection{Preliminaries: Backbone and Post-Training}
\label{sec:prelim}

\paragraph{Model}
\base~\cite{cosmos3} is a mixture-of-transformers model~\cite{liang2024mot} with two towers coupled by joint attention: a \emph{reasoner} tower that processes the text instruction and the conditioning tokens autoregressively, and a \emph{generator} tower that denoises video latents and action tokens with a rectified-flow objective~\cite{lipman2023flow,liu2023rectified}.
Video is tokenized by a causal spatio-temporal VAE~\cite{wan2025}.
Actions of every embodiment are expressed in a unified representation: the relative end-effector translation, a 6-D rotation~\cite{zhou2019rot6d}, and a gripper state per arm, zero-padded to a fixed width.
They enter and leave the token stream through embodiment-specific \emph{action heads}.
A head consists of an input projection, an output projection, and a modality embedding, with one row per embodiment.
When used as a policy, the model conditions on the current observation and the instruction and jointly generates the following video frames and the corresponding chunk of actions.
Our method uses only the action pathway of this backbone; the video pathway is kept as it is in the recipe we inherit.

\paragraph{Objective}
For a dataset $\calD$ of windows, each consisting of an observation, an instruction, the following video frames, and the corresponding actions, the training loss is
\begin{equation}
  \calL(\theta;\calD) = \lambda_v\,\calL_{\mathrm{video}}(\theta;\calD) + \lambda_a\,\calL_{\mathrm{action}}(\theta;\calD),
  \label{eq:loss}
\end{equation}
two flow-matching losses on the video latents and the action tokens, weighted by $\lambda_v$ and $\lambda_a$.
For a VLA, $\calL$ reduces to the action term.

\paragraph{Post-training recipe}
Given a target dataset $\calD_\calT$, post-training minimizes \cref{eq:loss} over a trainable subset $\theta$ of the weights, the generator tower together with its timestep embedding, its latent projections, and the action head of the target embodiment, while the reasoner tower and the tokenizers stay frozen.
The target embodiment is new, so its action-head row is initialized randomly.
We write $\PT_T(\theta_0;\calD_\calT)$ for the weights after $T$ optimizer steps of this recipe from the initialization $\theta_0$.
The standard pipeline uses the pre-trained weights, $\theta_0=\theta_{\mathrm{pre}}$.

\subsection{Few-Shot Post-Training}
\label{sec:problem}

We consider the regime in which $\calD_\calT$ holds only $n$ demonstrations per task, with $n$ small.
We deliberately keep the recipe fixed, so that the only free variable is the initialization $\theta_0$, and study which $\theta_0$ maximizes the closed-loop performance of $\PT_T(\theta_0;\calD_\calT)$ when $\calD_\calT$ is small.
\cref{fig:teaser} shows that $\theta_{\mathrm{pre}}$ is a poor choice.
Pre-training minimizes $\E_{e}[\calL(\theta;\calD_e)]$ over a mixture of embodiments $e$, \ie, it looks for one set of weights that is good for all of them at once.
Few-shot post-training needs something else: weights from which a short optimization on a \emph{new} embodiment lands in a good region.
The two coincide only if the per-embodiment solutions lie close to the joint solution.
We hypothesize that they do not, and that this mismatch, together with the randomly initialized action head, is what makes post-training from $\theta_{\mathrm{pre}}$ data hungry.

\subsection{Meta Pre-Training}
\label{sec:meta}

We therefore add a stage that directly optimizes $\theta$ for being a good starting point of $\PT$ on an unseen embodiment, using source embodiments as stand-ins for the target.
The procedure is summarized in \cref{alg:meta}.

\begin{algorithm}[t]
\caption{\Method}
\label{alg:meta}
\begin{algorithmic}[1]
\Require pre-trained weights $\theta_{\mathrm{pre}}$; source datasets $\{\calD_e\}_{e=1}^{E}$; fine-tuning step $\FT$; $K$, $w$, $k$, $N$; schedule $\epsilon_1\ge\dots\ge\epsilon_N$
\State $\theta\gets\theta_{\mathrm{pre}}$; initialize the shared action head randomly
\For{$i=1,\dots,N$}
  \State sample $e\sim\Unif\{1,\dots,E\}$ and $K$ demonstrations of $\calD_e$
  \State $\calS_e\gets$ $w$ windows from each demonstration
  \State $\tilde\theta_e\gets\FT_k(\theta;\calS_e)$ \Comment{$k$ fine-tuning steps on $\calS_e$, fresh optimizer state}
  \State $\theta\gets\theta+\epsilon_i\,(\tilde\theta_e-\theta)$
\EndFor
\State \Return $\thmeta=\theta$ \Comment{trainable set and shared action head}
\end{algorithmic}
\end{algorithm}

\paragraph{Source episodes}
We use $E$ robot datasets that are already part of the pre-training mixture, each recorded with a different embodiment (\cref{sec:setup} lists them); no new data is collected.
Every source is read through the same windowing and transformation pipeline as the target, so meta-training samples have exactly the format that post-training samples have.
An episode samples an embodiment $e\sim\Unif\{1,\dots,E\}$ and $K$ demonstrations of $e$, and takes $w$ windows from each of them.
We call the resulting set of $Kw$ windows the \emph{support set} $\calS_e$; it is the only data that the inner loop of the episode sees.
Sampling whole demonstrations rather than windows matters: the downstream data consists of a few complete trajectories, and the support set should have the same structure, a few trajectories with many windows each.

\paragraph{Shared action head}
The target embodiment arrives without an action head.
To meta-learn the head as well, all source embodiments share one spare action-head row during \method, initialized from scratch like the downstream head.
At post-training time this row is copied into the row of the target embodiment, so the target head starts from a meta-learned initialization instead of a random one.

\paragraph{Inner loop}
From the current weights $\theta$, the inner loop runs $k$ optimizer steps of ordinary fine-tuning on the support set $\calS_e$, \ie, on the windows sampled from the $K$ demonstrations of embodiment $e$, with the objective of \cref{eq:loss}; the optimizer state is reset at every episode.
We write $\tilde\theta_e=\FT_k(\theta;\calS_e)$ for the fine-tuned weights.
Each inner step covers the whole support set, so the $k$ steps make $k$ passes over the $K$ demonstrations, a miniature of post-training, which makes many passes over a few demonstrations.

\paragraph{Outer update}
The outer update interpolates between the current and the adapted weights,
\begin{equation}
  \theta\leftarrow\theta+\epsilon\,(\tilde\theta_e-\theta),
  \label{eq:outer}
\end{equation}
with $\epsilon$ annealed over the $N$ episodes.
This is the first-order update of~\cite{nichol2018reptile}.
In expectation, \cref{eq:outer} converges to the point at which $k$ steps of fine-tuning on a randomly drawn source embodiment no longer move the weights on average.
That point lies at the center of the per-embodiment $k$-step solutions: a few fine-tuning steps suffice to reach the solution of every source embodiment.
The source embodiments are stand-ins for the target.
The target embodiment is never seen during \method, but the property being optimized, that a few demonstrations and a few steps are enough to adapt, is enforced across embodiments that differ in robot, cameras, and action space, and we expect it to carry over to a new embodiment in the same way that a meta-learned initialization transfers to unseen tasks in few-shot learning~\cite{finn2017maml,nichol2018reptile}; \cref{sec:exp} confirms this empirically.
Nichol \etal~\cite{nichol2018reptile} further show that for $k>1$ the expected update contains, besides the average gradient, a term that increases the inner product between the gradients of different minibatches of the same task.
In our setting that term rewards weights for which fitting a few windows of an embodiment also reduces the loss on its remaining windows, which is precisely the generalization that few-shot post-training requires.

\paragraph{Why an interpolation update}
Three properties make the interpolation update the right tool at this scale.
\emph{(i) Long inner horizon.}
The downstream adaptation runs for thousands of optimizer steps.
Meta-gradient methods such as first-order MAML, whose update is the query-set gradient at the adapted point, are tied to the short inner horizon they are trained with; in preliminary experiments with a short inner loop they produced initializations whose offset from $\theta_{\mathrm{pre}}$ the subsequent post-training neither used nor undid.
The interpolation update has no query set, and its $k$ is limited only by compute.
\emph{(ii) Memory.}
\cref{eq:outer} needs one extra copy of the trainable parameters and nothing else; the model's own sharded parameters~\cite{zhao2023fsdp} are the fast weights, and the standard trainer, optimizer, activation checkpointing, and checkpoint code run the inner loop unchanged.
This is what makes meta-training the entire trainable set affordable on a few GPUs, whereas an inner loop with per-step fast weights or a query gradient at the adapted point would not fit.
\emph{(iii) Compatibility.}
The outer update is a convex combination of weights, so $\theta$ remains a valid checkpoint of the unchanged architecture at every iteration and is loaded by the post-training recipe without modification.
Meta-training the whole trainable set rather than only the heads is important: in preliminary experiments, meta-learning only the action heads (an ANIL-style variant~\cite{raghu2020anil}) gave at most marginal gains, because full-parameter post-training then overwrites the small meta-learned subset.

\subsection{Ranking Policy Learning}
\label{sec:rpl}

\paragraph{Visual bias of few-shot policies}
A good initialization does not remove a second difficulty of few-shot post-training.
With $n$ demonstrations per task, the training set contains a few trajectories per instruction, and in each of them the visual context, the objects and their layout at the start of the trajectory, is almost perfectly predictive of the behavior that follows.
The policy can therefore fit the demonstrations without reading the instruction, and the objective of \cref{eq:loss} provides no signal against this shortcut: it only asks the prediction under the correct instruction to be accurate.
At test time in a scene that is shared by several tasks, such a policy executes the behavior that the scene suggests rather than the one that was instructed, and in our rollouts this is the dominant failure mode on the suites whose tasks share one scene (\cref{sec:discussion}).
We refer to this effect as \emph{visual bias}.

\paragraph{Objective}
\Rpl (\RPL) adds a signal that depends on the instruction.
For each training window with observation $o$, instruction $c$, and target actions, we sample a negative instruction $c^-$ from the other tasks of the target set and evaluate the action flow-matching error under both instructions, with the same noised actions $x_t$, the same timestep $t$, and the same observation:
\begin{equation}
  E^{+} = \big\|v_\theta(x_t,t,o,c)-v^\star\big\|^2,\qquad
  E^{-} = \big\|v_\theta(x_t,t,o,c^-)-v^\star\big\|^2,
  \label{eq:rpl_errors}
\end{equation}
where $v_\theta$ is the predicted and $v^\star$ the target velocity of the action tokens.
The ranking loss
\begin{equation}
  \calL_{\mathrm{rank}} = \E\big[\max\big(0,\,m+E^{+}-E^{-}\big)\big]
  \label{eq:rpl}
\end{equation}
requires the prediction under the wrong instruction to be worse than under the correct one by a margin $m$, and the post-training objective becomes $\calL+\lambda_r\,\calL_{\mathrm{rank}}$.
The hinge is inactive once the margin is met, so windows whose actions already depend on the instruction are left alone, and it pushes the two predictions apart otherwise, so that the instruction, not only the scene, determines the action.
The negative branch decodes actions only and adds no video objective under the wrong instruction; windows whose instruction is dropped for classifier-free training contribute no ranking loss.
\RPL adds no parameters and no data, and costs one additional forward and backward pass of the generator per window.

\subsection{Post-Training with Meta Pre-Trained Weights}
\label{sec:downstream}

Post-training from $\thmeta$ differs from the baseline in two places: the trainable set is loaded from $\thmeta$ instead of $\theta_{\mathrm{pre}}$, and the target action-head row is initialized from the shared meta-learned head.
Where \RPL is used, the ranking loss of \cref{eq:rpl} is added to the objective.
The frozen reasoner tower and tokenizers, the data pipeline, the optimizer, the schedule, the number of steps, and the inference procedure are unchanged.
The same $\thmeta$ is used for every target suite and every data budget in \cref{sec:exp}; \method is a one-time cost that is amortized over all downstream targets.
